\chapter{Literature Survey}

\section{Foundations of Fact Verification}

\subsection{Automatic Fact Verification is a Text Classification Problem}
Text classification problems are defined as problems where you are given a textual input, and then you assign a label to said text \parencite{minaee2021deep}. Fact verification in most of its forms is a text classification problem where, given a claim and a piece of evidence both expressed as text, you determine if the claim is entailed or refuted, or also if there is insufficient information. This was how it was done in the standard baseline FEVER \parencite{thorne2018fever}, which formalised it as a 3-label task (SUPPORTED, REFUTED, NOT ENOUGH INFO). As it is a text classification task, you can apply the methods of text classification to provide a solution to automatic fact verification. 

\subsection{Implementing Fact Verification}
Most solutions for solving text classification problems require you to convert your input as a feature vector. Logistic regression \parencite{cox1958regression} is a simple method which uses this, where after converting the table and claim into a vector, using a technique such as TFIDF \parencite{salton1988term} which creates sparse vectors, you create a hyperplane to separate the classes of entailment and refutation in a high-dimensional feature space. 

\vspace{\baselineskip}
\noindent More complex solutions use dense feature vectors (embeddings) like in supervised learning with multi-layer perceptrons \parencite{christopher2006pattern}, where, you compute forward passes to arrive at logits which can be converted into probabilities for each text classification class, and then adjust the forward pass weights  with backpropagation based on the loss. This can be particularly effective when provided with contextual embeddings from pretrained transformers \parencite{vaswani2017attention} with encoder-only and decode-only backbones. For example, BERT \parencite{devlin2019bert}, an encoder, can more effectively capture semantic meaning in statements through attention as it is pretrained through masked token prediction and next sentence prediction. This is important as static word embeddings are the same for a word even in different contexts. Finetuning a classification head MLP on top of a contextual embedder like BERT for text classification tasks, can produce impressive results like in (citations?). Autoregressive decoder-only models such as GPT-2 \parencite{radford2019language} and Qwen3 \parencite{yang2025qwen3} are pretrained by next-token prediction, and are also adapted for use in text classification tasks, including fact verification such as in (citations?). Such approaches will be based on fine-tuning the sequence classification head on the final token (citations?), or prompting instruction-tuned models with zero- or few-shot prompts (citations?) and parsing the resulting label predictions (or label-token probabilities directly). The baseline FEVER \parencite{thorne2018fever} evaluated many of these methods paired with fetching relevant natural language statements from Wikipedia.


\subsection{Multi-hop Reasoning}
A key concept in fact-verification is multi-hop reasoning, which in this context is defined as verifying a claim by connecting multiple pieces of information through a sequence of linked inference steps across one or many sources rather than just a direct fact. For example, to verify the hypothetical claim “The winner of Wimbledon 2021 hates pineapples”, you have to get information about the winner of Wimbledon in 2021 and a separate piece of information on whether they do or do not hate pineapples (2-hop). HotpotQA \parencite{yang2018hotpotqa} describes multi-hop reasoning categories as involving bridge entities, where you use the entity from one source to find the other source and comparison reasoning, which is where the question is based on a shared entity. HoVer \parencite{jiang2020hover} built on this for multi-hop open-domain fact verification up to four Wikipedia article hops, with performance degrading significantly for state-of-the-art models as hops increased, showing a need for effective multi-hop reasoning in fact verification pipelines that retrieve information and predict. 

\section{Representing Tables}
testing12343

\subsection{Methods of Representing Tables}

\subsection{Applications of Table Representation for Single Table Classification Tasks} 

\section{Evidence Retrieval and Reasoning with Multiple Pieces of Evidence}

\subsection{Need for Evidence Retrieval}

\subsection{Methods of Evidence Retrieval}

\subsection{Multi-hop Reasoning with Multiple Pieces of Evidence}

\section{Bringing it all together: Multi-table Multi-hop Fact Verification}

\subsection{End-to-end Application and Evaluation: TRACE}

\subsection{Future Directions}